% Authoritative exercise source. Preserve all derivations and diagrams.

\ExerciseStart{G2.02}{Sommare e moltiplicare due ampiezze}{Guida esplicita}

\BookHeading{Problema, dati e obiettivo}

Siano \(z=1+i\) e \(w=1-i\). Calcolare \(z+w\), \(zw\) e \(z^2\). Interpretare somma e prodotto nel piano complesso e confrontare le lunghezze ottenute.

\ExTheory{\ExRef{c2-prodotto-complesso}{2.02}}{La somma si esegue componente per componente. Il prodotto si sviluppa distributivamente usando \(i^2=-1\). In forma polare \((r e^{i\theta})(s e^{i\phi})=rs e^{i(\theta+\phi)}\): i moduli si moltiplicano e le fasi si sommano.}

\ExPrediction{I due numeri sono coniugati: hanno uguale lunghezza e angoli opposti. Prevedere la direzione del loro prodotto e confrontarla con quella della loro somma.}

\ExFigure{\begin{tikzpicture}[>=Stealth,exlabel]
\draw[exaxis] (-0.35,0)--(1.55,0) node[right] {$\operatorname{Re}$};\draw[exaxis] (0,-1.35)--(0,1.55) node[above] {$\operatorname{Im}$};\draw[->,Accent,line width=1.1pt] (0,0)--(1,1) node[anchor=south west] {$ z $};\draw[->,Secondary,line width=1.1pt] (0,0)--(1,-1) node[anchor=north west] {$ w $};
\end{tikzpicture}}{Fattori e addendi hanno uguale modulo e fasi opposte.}


\begin{ExTasks}

\item Calcolare la somma.

\item Sviluppare tutti i termini del prodotto.

\item Ricavare i moduli e gli angoli dei fattori.

\item Confermare il prodotto in forma polare.

\item Calcolare il quadrato di z e confrontare le due trasformazioni.

\end{ExTasks}

\ExSolution

\ExStep{Sommare le componenti}

\ExEq{z+w&=(1+1)+i(1-1)=2.}
La costruzione punta-coda porta dal punto \((1,1)\) al punto \((2,0)\).

\ExStep{Svolgere il prodotto}

\ExEq{zw&=(1+i)(1-i)=1\cdot1+1(-i)+i\cdot1+i(-i)\\&=1-i+i-i^2=1+1=2.}

\ExStep{Passare alla forma polare}

\ExEq{|z|=|w|&=\sqrt{1^2+1^2}=\sqrt2,\\\arg z&=\pi/4,\quad\arg w=-\pi/4.}
I quadranti sono rispettivamente il primo e il quarto.

\ExStep{Interpretare il prodotto come trasformazione}

\ExEq{zw&=(\sqrt2 e^{i\pi/4})(\sqrt2 e^{-i\pi/4})
=2e^0=2.}
Moltiplicare z per w applica una rotazione di \(-45^\circ\) e una dilatazione di fattore \(\sqrt2\).

\ExStep{Confrontare con il quadrato}

\ExEq{z^2&=(1+i)(1+i)=1+i+i+i^2=2i,\\z^2&=(\sqrt2)^2e^{i(\pi/4+\pi/4)}=2e^{i\pi/2}=2i.}
Il modulo è ancora 2; la direzione è verticale perché la rotazione ora è \(+45^\circ\). L’uguaglianza numerica fra somma e prodotto nel primo confronto dipende dai dati scelti.

\par\noindent\begin{minipage}{\linewidth}
\ExFigure{\begin{tikzpicture}[>=Stealth,exlabel]
\draw[exaxis] (-0.35,0)--(2.55,0) node[right] {$\operatorname{Re}$};\draw[exaxis] (0,-0.35)--(0,2.55) node[above] {$\operatorname{Im}$};\draw[->,Accent,line width=1.1pt] (0,0)--(2,0) node[anchor=south west] {$ zw $};\draw[->,Secondary,line width=1.1pt] (0,0)--(0,2) node[anchor=south west] {$ z^2 $};
\end{tikzpicture}}{I due prodotti hanno modulo 2 e fasi rispettivamente 0 e \(\pi/2\).}


\ExInterpret{La somma combina contributi mediante traslazione dei segmenti; il prodotto trasforma un contributo mediante rotazione e dilatazione. Questa distinzione sarà necessaria per seguire interferenza e porte di fase.}

\ExLab{Nella sezione 2.02 impostare i due numeri nello strumento esplorativo. Confrontare somma e prodotto, poi sostituire w con z. Controllare che la costruzione geometrica e il calcolo distributivo descrivano lo stesso risultato.}


\end{minipage}
